\section{The AI learning assistant}
\label{sec:ai}

\subsection{What makes it more than a chatbot}

The assistant can \emph{act} on the user's library --- create a branch of posts, rewrite an article, move a
folder, attach a glossary term --- and is structurally incapable of doing any of it by itself.

Its entire output is one JSON object: no database handle, no file access, no tool-calling API, no way to
name a row id --- it addresses everything by title. Proposed operations are parsed on the device, validated
against the live database, rendered as a checklist, and applied only on an explicit tap inside one Room
transaction. Figure~\ref{fig:aiflow} is that write path on one page.

% The AI write path: nothing reaches Room until the user approves it.
\begin{figure}[H]
\centering
% Node widths are fixed in cm, so scale the finished diagram to the measure rather
% than re-laying it out whenever the body font changes.
\resizebox{\linewidth}{!}{%
\begin{tikzpicture}[
  font=\scriptsize,
  x=1cm, y=1cm,
  step/.style={draw=mochaBrown, fill=mochaCream, rounded corners=2pt, align=center,
               minimum height=9mm, text width=#1, inner sep=3pt, anchor=center},
  gate/.style={draw=mochaError, line width=0.7pt, fill=mochaCreamDark, rounded corners=2pt,
               align=center, minimum height=9mm, text width=#1, inner sep=3pt, anchor=center},
  ext/.style={draw=mochaSage, fill=mochaCreamDark, rounded corners=2pt, align=center,
              minimum height=9mm, text width=#1, inner sep=3pt, anchor=center},
  tag/.style={font=\scriptsize\itshape\color{mochaDeep}, align=center, anchor=center},
  go/.style={-{Stealth[length=1.9mm]}, draw=mochaDeep},
  loop/.style={-{Stealth[length=1.9mm]}, draw=mochaSage, dashed},
]

% --- request row -----------------------------------------------------------
\node[step=2.1cm] (msg)  at (-6.4, 0)    {user message\\+ mode};
\node[step=2.4cm] (idx)  at (-3.3, 0)    {build kbIndex\\\itshape titles only, $\le$4\,KB};
\node[ext=2.1cm]  (gw)   at (-0.2, 0)    {gateway\\$\rightarrow$ DeepSeek};
\node[step=2.5cm] (env)  at ( 3.1, 0)    {parse one JSON\\envelope};

% --- the three outcomes ----------------------------------------------------
\node[step=2.3cm] (answer) at (-6.4,-2.6) {answer bubble\\\itshape streamed live};
\node[ext=2.5cm]  (tools)  at (-2.6,-2.6) {ContextTools\\\itshape 7 read-only queries};
\node[step=2.8cm] (val)    at ( 3.1,-2.6) {ActionValidator\\\itshape schema, refs, parents,\\\itshape cycles, caps};

\node[gate=3.0cm] (review) at ( 3.1,-5.0) {\bfseries Review screen\\\itshape per-action checkbox,\\\itshape preview, location};
\node[step=3.0cm] (tx)     at ( 3.1,-7.1) {ActionExecutor\\\itshape one Room transaction};
\node[step=2.3cm] (undo)   at (-0.9,-7.1) {Undo snapshot};

% --- arrow labels, placed explicitly so nothing collides -------------------
\node[tag] at (4.35,-1.25) {type=actions};
\node[tag] at (4.05,-3.75) {valid};
\node[tag] at (1.85,-6.05) {approved\\subset};
\node[tag] at (-5.35,-0.85) {type=answer};
\node[tag] at (0.55,-1.85) {type=\\context\_request};
\node[tag] at (-3.40,-1.55) {results,\\$\le$4 rounds};

\draw[go] (msg) -- (idx);
\draw[go] (idx) -- (gw);
\draw[go] (gw)  -- (env);
\draw[go] (env.south) -- (val.north);
\draw[go] (val.south) -- (review.north);
\draw[go] (review.south) -- (tx.north);
\draw[go] (tx.west) -- (undo.east);
\draw[go] (env.south west) to[out=215, in=25] (answer.north east);
\draw[loop] (env.south west) to[out=250, in=15] (tools.east);
\draw[loop] (tools.north) to[out=100, in=250] (idx.south);

\node[align=left, text width=3.5cm, anchor=west, font=\scriptsize\color{mochaMuted}]
      at (5.0,-2.6) {Any invalid action rejects the \emph{whole} batch with a readable error list --- a
      half-applied proposal would be worse than none.};
\node[align=left, text width=3.5cm, anchor=west, font=\scriptsize\color{mochaError}]
      at (5.0,-5.0) {The only path from the model to the database. Destructive actions arrive unchecked
      and need a confirmation.};
\end{tikzpicture}%
}
\caption{The AI write path. The model never sees the database and never writes to it: it emits a JSON
envelope, the device validates it against live data, the user approves individual actions, and one Room
transaction applies them all-or-nothing.}
\label{fig:aiflow}
\end{figure}


\subsection{Four modes}

The mode travels with every request as a plain string; anything outside the set is rejected with HTTP 400.
It selects a short brief appended to the system prompt.

\noindent
\begin{tabularx}{\linewidth}{L{2.2cm} X}
\toprule
\textbf{Mode} & \textbf{The brief appended to the system prompt} \\
\midrule
\textbf{Answer} &
Answer conversationally, propose nothing, use \code{context\_request} to look things up. \\
\addlinespace
\textbf{Suggest} &
Recommend posts, paths, links, resources or terms as a batch; ``prefer small, high-value suggestions''. \\
\addlinespace
\textbf{Modify} &
Propose concrete changes --- create, move, edit, link, tag --- with ``genuinely educational markdown''.
The mode that generates whole articles and learning paths. \\
\addlinespace
\textbf{Organize} &
Analyse a branch and propose restructuring: moves, new folders, links, duplicate detection. \\
\bottomrule
\end{tabularx}

\medskip
The action-protocol block goes out in \emph{all four} modes; only the brief differs, so the Answer-mode
guard lives in \code{ChatRepository} rather than the prompt --- a prompt instruction is a request, not a
constraint. Actions proposed in Answer mode never become a reviewable row; the summary is shown as prose
plus \emph{``switch to Suggest or Modify to review them''}.

\subsection{The gateway}

The app never talks to DeepSeek. It talks to a stateless Express process of about 260 lines that holds the
key, builds the request, normalises errors and relays the stream. It has no database: every request carries
its own history. Three routes:

\begin{itemize}
  \item \code{GET /v1/health} $\rightarrow$ \code{\{ok, model, streaming\}}: \code{ok} is whether the key is
        configured, so a keyless gateway reports itself unusable up front, and \code{streaming} saves a
        probe for a 404.
  \item \code{POST /v1/chat} --- buffered. Body \code{\{mode, messages[], kbIndex, toolResults?\}},
        response \code{\{reply, usage\}}.
  \item \code{POST /v1/chat/stream} --- same body, answered as Server-Sent Events.
\end{itemize}

Upstream it always asks for \code{response\_format: json\_object} and \code{max\_tokens} of 8192, because
DeepSeek otherwise stops at its default length and returns a truncated envelope.

\begin{lstlisting}[language=bash, caption={A request, its reply, and a normalised error. The whole answer is
the \texttt{reply} string, itself a JSON envelope parsed on the device.}]
POST /v1/chat
{ "mode": "modify",
  "messages": [{"role":"user","content":"Give me a path into consensus algorithms"}],
  "kbIndex": "- Distributed Systems (branch)\n  - Paxos (post) [reading]\n",
  "toolResults": null }

200 OK
{ "reply": "{\"type\":\"actions\",\"summary\":\"Create a Consensus path\",...}",
  "usage": {"prompt_tokens": 1841, "completion_tokens": 962} }

502
{ "error": "The reply was cut off before it finished. Try asking for less at once.",
  "retryable": true }
\end{lstlisting}

The key is read from \code{DEEPSEEK\_API\_KEY} in the gateway's \code{.env}, in neither git nor the build;
the APK holds only the gateway address, an editable Settings field.

Every failure reaches the app as \code{\{error, retryable\}}: missing key $\rightarrow$ 503 non-retryable;
unknown mode or empty \code{messages} $\rightarrow$ 400 non-retryable; upstream 429 $\rightarrow$ 429; other
upstream failures $\rightarrow$ 502 carrying DeepSeek's message; timeout $\rightarrow$ 504 (60\,s buffered,
120\,s streaming); \code{finish\_reason == "length"} or a blank reply $\rightarrow$ 502 retryable. Streaming
errors arrive as a frame, headers being already flushed, and a stream ending without \code{done} fails.

\subsection{Choosing what the model may see}

Sending the whole knowledge base would be expensive, would break as the library grew, and would ship every
private note off the device on every message. The model gets a map and must ask for the territory.

\code{KbIndex.build} walks the tree depth-first in sibling order, one indented line per node: title, type,
and reading status where it is not \code{NONE}. No bodies, tags, ids or timestamps, and a hard cap of
\code{MAX\_CHARS = 4000} characters, truncated mid-line.

\begin{lstlisting}[language=bash, caption={The compact KB index --- the only context sent unprompted.}]
- Distributed Systems (branch)
  - Consensus (folder)
    - Paxos (post) [reading]
    - Raft (post) [in_progress]
\end{lstlisting}

If it needs more it replies with a \code{context\_request}; the device runs the queries locally and returns
the results. Seven read-only operations exist in \code{ContextTools}:

\noindent
\begin{tabularx}{\linewidth}{L{3.2cm} L{2.1cm} X}
\toprule
\textbf{Operation} & \textbf{Argument} & \textbf{What comes back} \\
\midrule
\code{search\_posts} & \code{query} &
Up to 10 hits as \code{KIND: Title}, from the same full-text search the Search tab uses. \\
\code{get\_post} & \code{title} &
Title, tag names and the full Markdown body. The only operation that returns article text. \\
\code{list\_children} & \code{parentTitle} &
That container's children as \code{type: title}; blank returns the root branches. \\
\code{get\_backlinks} & \code{title} & Titles of posts whose bodies link to this one. \\
\code{search\_dictionary} & \code{query} &
Up to 15 glossary rows as \code{term: definition}; wildcards are stripped from the fragment. \\
\code{get\_tags} & --- & Every tag name, one per line. \\
\code{get\_related} & \code{title} &
Up to five from \code{PostRepository.related} (shared tags and link targets, no backlinks). \\
\bottomrule
\end{tabularx}

\medskip
Three caps apply: \textbf{8 queries} per request, each result truncated to \textbf{2\,000 characters}, and
\code{repeat(4)} in \code{ChatRepository} --- \textbf{four model calls per user message}, at most three of
them context requests. A \code{context\_request} on the fourth is refused with a retryable \emph{``Could not
gather enough context''} rather than looping.

Results come back as an ordinary \code{user} turn labelled \emph{``CONTEXT TOOL RESULTS \dots (read-only,
produced on-device)''} --- not the OpenAI \code{role: "tool"} turn, which DeepSeek rejects with HTTP 400
unless it carries a \code{tool\_call\_id}. Each reply carries \emph{``Shared with AI: $n$ notes''}, $n$
counting queries run.

\subsection{The action protocol}

The reply is one JSON object and nothing else, of three envelope types: \code{answer} carries \code{text}
and renders as a bubble; \code{context\_request} carries \code{queries} and triggers the round trip above;
\code{actions} carries a \code{summary} and the action list for the review screen.

\code{ActionParser} is forgiving about packaging and strict about content: it strips a code fence, takes the
substring from the first \code{\{} to the last \code{\}}, and Gson-parses it. A blank or missing
\code{type} is an error; prose that is no envelope is wrapped as an \code{answer} by \code{parseOrAnswer},
so forgetting the format is never fatal.

\begin{lstlisting}[language=bash, caption={An \texttt{actions} envelope; \texttt{ref} names an item created
earlier in the same batch.}]
{ "type": "actions",
  "summary": "Create a Consensus path with 4 posts",
  "actions": [
    {"op":"create_branch","title":"Distributed Systems","ref":"b1"},
    {"op":"create_folder","parentRef":"b1","title":"Consensus","ref":"f1"},
    {"op":"create_post","parentRef":"f1","title":"Raft","ref":"p1",
     "content":"Raft is a consensus algorithm ... see [[Paxos]].",
     "tags":["consensus"],"status":"READING"},
    {"op":"create_link","fromRef":"p1","toTitle":"Paxos"},
    {"op":"add_dictionary_entry","postRef":"p1","term":"quorum",
     "definition":"A majority of nodes.","meaningVi":"..."} ] }
\end{lstlisting}

Fourteen operations are implemented, listed below as \code{ActionValidator.knownOps}; \code{backend/prompts.js}
documents the same set, each side commented that changing one means changing the other.

\begin{xltabular}{\linewidth}{L{3.75cm} L{3.0cm} X}
\toprule
\textbf{Operation} & \textbf{Key fields} & \textbf{Effect when applied} \\
\midrule
\endfirsthead
\toprule
\textbf{Operation} & \textbf{Key fields} & \textbf{Effect when applied} \\
\midrule
\endhead
\code{create\_branch} & \code{title}, \code{ref} & New root container. \\
\code{create\_folder} & \code{parentRef} or \code{parentTitle}, \code{title}, \code{ref} &
New container under an existing or in-batch parent. \\
\code{create\_post} & + \code{content}, \code{tags}, \code{status} &
New article; tags created on demand, status defaults to \code{READING}. \\
\code{update\_post} & \code{postTitle}, \code{title}, \code{content}, \code{tags}, \code{status} &
Rewrites body, title, tag set and/or status. \\
\code{move\_post} & \code{postTitle}, \code{newParentTitle} & Re-parents a post. \\
\code{delete\_post} & \code{postTitle} & Deletes the post and everything hanging off it. \\
\code{create\_link} & \code{fromRef} or \code{fromTitle}, \code{toTitle} &
Appends \code{[[Target]]} to the source body, which creates the edge. \\
\code{remove\_link} & \code{fromTitle}, \code{toTitle} & Removes that wiki-link from the body. \\
\code{add\_tag}, \code{remove\_tag} & \code{postTitle}, \code{tag} & One tag on one post. \\
\code{set\_status} & \code{postTitle}, \code{status} &
\code{READING}, \code{IN\_PROGRESS} or \code{FINISHED}. \\
\code{set\_favorite} & \code{postTitle}, \code{favorite} & Stars or unstars a post. \\
\code{add\_resource} & \code{postTitle}, \code{type}, \code{title}, \code{url} &
Attaches a \code{YOUTUBE}, \code{ARTICLE}, \code{BOOK} or \code{OTHER} reference. \\
\code{add\_dictionary\_entry} & \code{postTitle} (optional), \code{term}, \code{definition},
\code{meaningVi} & Adds or updates a glossary term, post-scoped or global. \\
\bottomrule
\end{xltabular}

\medskip
Nothing here is vendor-specific: no function-calling schema, no tool role, no provider SDK --- a documented
JSON contract over ordinary chat completions. Swapping providers is an edit to \code{DEEPSEEK\_URL},
\code{MODEL} and the key in \code{.env}; no Android changes.

\subsection{Validation, review and undo}

\code{ActionValidator} takes the parsed actions and a live snapshot of every node, returning errors prefixed
\code{Action N:}. Its \emph{scope} --- existing nodes by lower-cased title, plus growing sets of what
earlier actions will create --- is how in-batch references work. An action may declare a \code{ref} that
later ones point at with \code{parentRef}, \code{postRef} or \code{fromRef}; it must \emph{already} be
registered, and registration happens \emph{after} the parent check, so a forward reference fails rather than
half-working. Titles resolve the same way, so a \code{set\_status} on a post \code{create\_post} proposes
three rows earlier validates.

It rejects: a duplicate \code{ref} or an \code{op} outside the fourteen; more than 40 actions, a title over
200 characters, content over 50\,000, a term over 80, more than 20 tags; a \code{create} whose title exists
(live or in-batch) and an \code{update\_post} retitling onto a taken name, which must fail here rather than
inside the executor's transaction; a parent that does not exist or is a post; a \code{move\_post} that would
cycle, checked with the \code{TreeRules} helper drag-and-drop uses; a status outside the enum, an unknown
resource type, and missing \code{url}, \code{tag}, \code{term} or \code{definition}. A \code{create\_link}
whose \code{toTitle} is absent is deliberately allowed: a link to a post not yet written is legitimate, and
the reader renders it unresolved.

An invalid action is not silently dropped: the validator runs over the ticked subset, re-runs on every tick,
edit and relocation, and Apply stays disabled while any error stands --- unticking the offending row lets
the rest through. \code{ReviewChangesViewModel} rewrites each \code{Action N:} prefix into that action's
description and hangs it under the row. \code{apply()} re-validates before writing, so a stale batch cannot
slip past. The screen holds the batch in memory while the user reshapes it:

\begin{itemize}
  \item \textbf{Per-action selection.} Every row has a checkbox and an ``$n$ of $m$ selected'' counter,
        ticked except \code{delete\_post}, which cascades and cannot be undone.
  \item \textbf{Structure.} Rows are indented by depth \emph{inside the batch}: \code{ActionLabels.indent}
        walks \code{parentRef} and \code{parentTitle} chains with a cycle guard, anchoring at depth 0 on an
        existing parent, so a path renders as the tree it will become.
  \item \textbf{Previews and diffs.} A \code{create\_post} row previews its Markdown through the reader's
        Markwon instance; an \code{update\_post} row a before/after from \code{TextDiff} --- honestly
        primitive, a whole-document comparison, not a line diff: right for a rewrite, poor for one word.
  \item \textbf{Changing the destination.} Creatable and movable rows offer a picker of every branch and
        folder by full path (\emph{Distributed Systems / Consensus}) plus ``top level'', rewriting
        \code{parentTitle} and clearing \code{parentRef}.
  \item \textbf{Editing before saving.} Title and body can be rewritten in place; a retitle rewrites every
        other action addressing the old title, so tags, links and glossary entries follow the rename.
  \item \textbf{Confirmation for destruction.} A ticked destructive row raises a banner, and Apply a dialog
        naming each doomed post and its tags, links, resources and glossary entries, all beyond Undo.
\end{itemize}

Apply hands the ticked subset to \code{ActionExecutor}, which resolves refs to ids and runs every operation
in one \code{db.withTransaction}. It is all-or-nothing: if operation seven throws, Room rolls back one to
six and reports \emph{``Could not apply these changes. Nothing was written.''} Content-backed
\code{posts\_fts} needs no work here (\S\ref{sec:data}).

Undo is real but narrow: \textbf{in memory} (one \code{@Volatile} \code{UndoSnapshot} on
\code{ChatRepository}), \textbf{per batch and single level} (a second Apply overwrites it, process death
loses it), and offered from \textbf{one place}, the snackbar shown after Apply. Created nodes are deleted in
reverse order, so children go before parents; nodes the batch merely touched are restored from a
before-image captured inside the transaction --- title, content, status, tag set, favourite flag, and
parent, moved back only if it changed --- taken once per node on first touch, so several edits to one post
still restore its pre-batch state. Resources are removed by id; glossary terms need two cases because
\code{addTerm} upserts, so one the batch created is deleted and one it overwrote restored.

\code{delete\_post} is deliberately \textbf{not} undoable, and the snapshot's header comment says so:
deletion cascades to children, links, tags, resources and glossary rows, and resurrecting that sub-graph is
a feature in its own right. The app says so at every step, ending with \emph{``Undone --- but deleted posts
stay deleted.''}

\subsection{Streaming}

When health advertises \code{streaming}, replies arrive over SSE, and an older gateway still works on the
buffered endpoint. The transport is raw OkHttp: Retrofit's Gson converter buffers the whole body and leaves
nothing to stream. \code{SseChatClient} exposes a cold \code{Flow<StreamFrame>} and, since a blocked socket
read ignores cancellation, registers \code{invokeOnCompletion \{ call.cancel() \}} so leaving the screen
closes the socket; its 130\,s read timeout outlasts the gateway's 120\,s budget, which measures the gap
between chunks.

\code{SseFrames} is a framework-free line decoder, unit-tested on the JVM. It ignores blank lines, SSE
comments (the gateway writes \code{: open} to flush headers early) and the \code{[DONE]} sentinel, and
decodes \code{\{delta\}}, \code{\{done, reply\}} and \code{\{error, retryable\}}, the last two terminal. The
whole reply is repeated in the \code{done} frame, so the envelope is parsed from one authoritative string,
not re-joined deltas. A stream dying \emph{before} the first token falls back to the buffered endpoint once;
one dying later surfaces the error, since a retry would duplicate a half-rendered bubble.

Partial rendering is awkward because the wire carries JSON, not prose: the visible text is the value of
\code{"text"} inside \code{\{"type":"answer","text":\dots\}}, and the buffer is not valid JSON until the
final brace. Waiting would defeat streaming, so \code{StreamingAnswerExtractor} (plain Java) scans the
buffer and classifies it as \code{answer}, \code{actions}, \code{context\_request}, \code{unknown} (JSON so
far, \code{type} not yet seen) or \code{prose} (no opening brace). For an answer it decodes as much of
\code{text} as arrived, unescaping \code{\textbackslash n}, quotes and \code{\textbackslash uXXXX} and
stopping cleanly on an escape split across chunks; for the other two it returns nothing, so raw protocol
JSON is never shown and the bubble reads \emph{``Mocha is working\dots''} The streaming bubble is never
written to \code{chat\_messages}, so a dead stream leaves no half-finished row.

\subsection{Privacy}

The assistant contradicts the local-first commitment of \S\ref{sec:intro}, so the boundary is drawn
explicitly. \textbf{What leaves the device}, only on sending a message in the AI tab: the mode string; the
session's
history (user turns and completed assistant turns --- pending batches, error bubbles and other sessions are
filtered out); the KB index; and the context-query results. Article text leaves only through
\code{get\_post}, for a post the model named.

\textbf{What never leaves}: the Room database file, any post the model did not ask for, tags, backlinks,
resources, settings, backups, and the other four tabs. No analytics SDK, crash reporter or image loader is
in the build (\S\ref{sec:tech}); \textbf{the key} lives in the gateway's \code{.env} alone, the APK holds no
credential, and Auto Backup is disabled.

One honest weakness: the device-to-gateway hop is cleartext HTTP, the gateway being a local process with no
certificate. It is declared in \code{network\_security\_config.xml}, user-installed CAs are not trusted, and
the DeepSeek hop is HTTPS; on a shared LAN the first hop is readable, and production would need TLS.

\textbf{When the gateway is unreachable} --- for most graders, most of the time --- every send begins with a
health ping, and failure degrades to the banner and retry of \S\ref{sec:features}. Nothing is queued or
lost, and no other feature touches the network.

\subsection{Where the code differs from the plan}

\code{docs/plan.md} predates the subsystem; several decisions changed on contact with DeepSeek.
\textbf{\S11} appends tool results as a \code{role: "tool"} message, calls the index ``titles/ids''
(\code{KbIndex} emits no ids) and forces an answer after three rounds where the code returns a retryable
error. \textbf{\S12} calls streaming optional; it shipped. \textbf{\S13} rejects the whole batch on any
validation failure, where the validator checks the ticked subset, and \textbf{\S14} rebuilds FTS rows in the
transaction, which content-backed \code{posts\_fts} makes unnecessary. \textbf{\S15} specifies a
\code{ReviewChangesActivity}; review is a Fragment destination in the single-activity app. \textbf{\S22}'s
``Shared with AI: 3 posts'' chip counts context queries, not distinct posts, and is in-memory only.

One layering shortcut: \code{ContextTools} reaches \code{nodeDao} and \code{knowledgeDao} directly ---
read-only, but the one place in the AI package that skips a repository (\S\ref{sec:self}).
